\chapter{文献分层与继承边界}
\label{ch:lit}

本章要判定的不是文献是否权威，而是模型中的每一层方程从哪里来、哪一层明确不从那里来、以及那一层在本文哪一处被使用。继承表示后文直接使用其方程结构。改写表示边界、状态或判据换成了热加工决策。不继承表示该文献的结论不能被写成本文的结果。条目只讨论可以核到作者、题目和年份的来源。不转写原文句子，不补卷期里没有逐条核对过的页码。

未在本文检索中出现的对象，只记为这次没有读到。未读到不等于不存在。肉类水分等温线在已读入的公开拟合表里条数为零，这是该表的范围，不是对全部肉类测量的否定。

\section{多孔介质与食品热湿传递}

\paragraph{要判定的量}
能量、动量和组分方程在体积平均之后还保留哪些项，以及蒸烤边界上哪些项必须改写。要判定的是冷点能否离开几何中心，不是某一烤箱的升温曲线能否被复现。

\paragraph{继承的文献层与不继承的层}
Whitaker 给出干燥过程中多孔介质里热量、质量和动量同时传递的体积平均理论\cite{whitaker1977}。本文继承平均之后能量、动量与组分方程的结构：固相骨架、液相和气相分别平均，界面交换写成源项，而不是在孔隙尺度上跟踪每一个弯月面。不继承干燥工艺的具体边界，不继承该文为木材或颗粒床选定的对流通量闭合。

Datta 把这一骨架用于食品加工中的同时热质传递\cite{datta2007}。本文继承这个用法：食品被看成含液相与气相的多孔介质，蒸发和冷凝是界面源，而不是把“蒸汽”写成一种更大的导热系数。不继承某一烘箱的实验拟合，也不把该文的某一组物料参数当作本文的主结果。

挤压与油炸多孔介质的系列工作给出蒸发—冷凝源项和压力驱动汽流。本文只把这两类项改写到蒸烤界面：冷凝时潜热进入表面，蒸发时潜热离开表面；压力若升高，只作为气相能否排出的定性约束。不继承油炸外部的热油流动，不继承挤压机的螺杆边界。

Feyissa、Gernaey 与 Adler-Nissen 给出对流烤箱中肉块的三维耦合场，并表明冷点可以偏离几何中心\cite{feyissa2013}。本文继承现象：边界不对称或内部导热不对称时，最低温度不必落在厚度中点。方程改为一维主算例加一块显式二维截面。截面不是该文的三维耦合场。不继承该文的网格、时间步和肉块物料数值。

\paragraph{方程、算法或协议}
体积平均后的能量方程保留四类贡献，写成规格中的结构而不是干燥器的边界条件：
\begin{align*}
\langle\rho c\rangle\partial_t T
&=-\nabla\cdot(q_{\mathrm{cond}}+q_{\mathrm{conv}})\\
&\quad+\dot m_\Gamma\Delta h+q_{\mathrm{rxn}}.
\end{align*}
$q_{\mathrm{cond}}$ 是有效导热，$q_{\mathrm{conv}}$ 是液相或气相携带的显热，$\dot m_\Gamma\Delta h$ 是界面相变，$\ q_{\mathrm{rxn}}$ 是反应热。动量平均只保留到足以说明气相能否在压差下离开表面，不在在线模型里求解完整的孔隙 Navier–Stokes 方程。组分方程保留液相分数与气相分数的交换。

蒸比干热更快进入内部，写成冷凝潜热为正的源，不写成“蒸汽导热更快”。干壳写成近表面渗透率下降，以及有效导热率下降，从而内部热阻上升。评分核没有求解这组平均方程。它用十六条控制体的导热和一条经验潜热通量代替。经验通量 $18000\,u_{\mathrm{st}}(1-a_w)$ 的量级是瓦每平方米，不是 $\Delta h$ 乘以由分压算出的质量流。模块里写着的 $2.3\times 10^{6}\,\mathrm{J/kg}$ 没有乘进战役步进。因此日志里的温度不是 Whitaker 方程的解。

冷点偏离几何中心在评分核里只有两个来源。\texttt{hetero}=1 把初始温度铺成 $\pm 4^\circ\mathrm{C}$ 的空间倾斜，最低温可以不在中心，但这是初始条件，不是边界驱动的偏移。\texttt{hetero}=2 把下标为 $5$ 的控制体导热降到 $0.45$ 倍，热阻位置固定在大约三分之一厚度处。这是一维夹层，不是\cite{feyissa2013}的三维几何。另有一块 24 乘 16 的显式截面：脂肪列在下标 8，20 min 后最低温度仍在绝热侧下标 23，偏离中心 \(13.75\,\mathrm{mm}\)。无脂肪时下标相同。所以这块截面复现的是“冷点可以不在中心”，没有复现“脂肪把冷点移到夹层背侧”。

\subsection{来源对照}
\begin{table}[ht]
\centering
\caption{传递骨架。方式一列只说明本文用到哪一层。}
\begin{tabularx}{\textwidth}{lXll}
\toprule
来源 & 使用的结果 & 方式 & 不继承 \\
\midrule
Whitaker, 1977 & 体积平均后的能量、动量与组分 & 结构 & 干燥工艺的边界 \\
Datta, 2007 & 食品中同时热质传递 & 用法 & 某一烘箱的拟合 \\
挤压与油炸系列 & 蒸发—冷凝源、压力驱动汽流 & 改到蒸烤界面 & 油炸外部流动 \\
Feyissa 等, 2013 & 冷点可偏离几何中心 & 现象 & 该文网格与物料数 \\
\bottomrule
\end{tabularx}
\end{table}
挤压与油炸系列在这次检索里没有单独的文献条目，表中只保留现象层的改写，不补未核对的卷期。

\paragraph{本节不能推出的结论}
引用体积平均理论不能推出本文已经求解多孔介质真值。一维夹层的冷点移动不能代替三维烤箱里的冷点云图。油炸文献中的外部流动不能写成蒸烤腔的风场。

\section{水分活度、收缩与各向异性}

\paragraph{要判定的量}
含水率 $M$ 与水分活度 $a_w$ 是否被写成同一个量，以及导热、厚度、渗透率里哪些参数有文献区间、哪些只能做灵敏度。要判定的还包括：已读入的等温线能约束哪一类物料、约束到哪一个温度。

\paragraph{继承的文献层与不继承的层}
食品水分等温线的综述传统给出 $a_w$ 与含水率的非线性关系。GAB 形式
\[
M=\frac{M_0 C K a_w}{(1-K a_w)\,(1-K a_w+C K a_w)}
\]
作为该关系的一种闭合被继承\cite{gab1980,sorption2026}。Timmermann 对多层吸附参数的比较说明，BET 与 GAB 的参数相关，数值上可以病态\cite{gab1980}。本文继承这一警告：三个参数不单独解释，相关结构必须保留。不继承未经测定的 $M_0,C,K$ 数值，也不把拟合表里的数写成新的实验测定。

Clausius–Clapeyron 关系使吸附热随温度移动等温线。本文继承“温度升高时单层容量下降”这一方向。不继承把活度当成含水率本身，也不继承测定温度以外的外推数值。

Choi 与 Okos 给出组分热物性随温度和组成变化的加和传统\cite{choi1986}。本文继承 $k$ 随水分、脂肪和温度变化的趋势。不继承把各向同性标量当成全部导热信息。各向异性测量给出 $k_\parallel\neq k_\perp$。本文把它写成张量结构。不把某一品种的单点测量当成普遍常数。

蛋白受热收缩与孔隙变化的工艺文献给出厚度和渗透率随蛋白进度变化。本文写成收缩方程。不把收缩定义成品。收缩是几何与阻力的演化，品质进度仍由动力学积分定义。

\paragraph{方程、算法或协议}
活度进入表面蒸发的驱动力是 $a_w p_{\mathrm{sat}}(T_s)-p_v$，不是 $M$ 本身。因此同一含水率在淀粉与蛋白上可以对应不同的界面蒸汽压。评分核没有反演 GAB 代数。它用单调代理
\[
a_w=1-\mathrm{e}^{-m/s}(1-0.15\,k),
\]
土豆、面团、肉糜的尺度 $s$ 分别取 $0.055$、$0.08$、$0.10$，$k$ 取 $0.80$、$0.75$、$0.70$，再裁剪到 $[0.02,0.99]$。土豆的尺度靠近已读表里 $60^\circ\mathrm{C}$ 附近的单层容量量级，但代理不是该表的拟合残差。肉糜的尺度是宽而弱的信息，因为该表中肉类曲线条数为 $0$。未读到不等于不存在：本文没有读到可引用的肉类等温线文件，所以肉类活度保持为开放参数，不用淀粉的 $M_0,C,K$ 冒充。

温度移动只在测定范围内解释。公开表中土豆的清洁曲线到 $60^\circ\mathrm{C}$ 为止。烤制表面高于此温度时，GAB 数值是外推，不能进入主结果。$M_0,C,K$、$k_\parallel/k_\perp$ 和收缩系数在没有文献区间时只进入灵敏度扫描。

各向异性写成
\[
\vect{k}_{\mathrm{eff}}=k_\perp(M,S)\,\vect{I}+(k_\parallel-k_\perp)\,\vect{e}\otimes\vect{e}.
\]
脂肪夹层的热阻是 $L_{\mathrm{fat}}/(k_{\mathrm{fat}}A)$，加在冷点路径上。评分核没有这个张量。均匀导热率是 $0.68\,\mathrm{W/(m\cdot K)}$，只有下标 $5$ 的控制体在 \texttt{hetero}=2 时降到 $0.45$ 倍，即 $0.306\,\mathrm{W/(m\cdot K)}$。这是标量缺损，不是各向异性。

\subsection{来源对照}
\begin{table}[ht]
\centering
\caption{界面与结构闭合。没有单独条目的综述只记传统，不补页码。}
\begin{tabularx}{\textwidth}{lXll}
\toprule
来源 & 使用的结果 & 方式 & 不继承 \\
\midrule
GAB 形式 \cite{gab1980,sorption2026} & \(a_w\) 与含水的非线性 & 形式 & 未测定的 \(M_0,C,K\) \\
Clausius--Clapeyron 传统 & 等温线随温度移动 & 方向 & 把活度当成含水 \\
组分加和 \cite{choi1986} & \(k\) 随水分和脂肪变化 & 趋势 & 各向同性即全部 \\
各向异性测量传统 & \(k_\parallel\neq k_\perp\) & 张量结构 & 单点当普遍常数 \\
收缩工艺传统 & \(L\) 与渗透率随进度变化 & 演化方程 & 收缩即品质 \\
\bottomrule
\end{tabularx}
\end{table}

收缩规格是
\[
\dot L=-\beta_s L_0\dot\alpha_p,\qquad K=K_0(L/L_0)^{n_K}\varepsilon_g^{n_\varepsilon}.
\]
评分核的厚度在一次运行内不变。结壳变量 $S$ 只进入热阻代理 $r_{\mathrm{proxy}}=\mathrm{clip}(\Delta T/8+1.2+2.5 S,\,0.8,\,8)$，不改变网格间距。渗透率不在状态里。因此日志不能检验收缩与内部蒸汽排出的耦合。

\paragraph{本节不能推出的结论}
清洁拟合的分位数不能描述全部食品，也不能验证冷点、探针或烤制。代理活度下降不能写成 GAB 三参数已被辨识。一维导热缺损不能写成各向异性测量。

\section{反应进度与热力当量}

\paragraph{要判定的量}
品质是一条温度阈值，还是一条速率方程的积分。要分开的是冷点上的热力进度、表面褐变和含水开关。颜色能不能代替中心进度，是本节必须否定的用法。

\paragraph{继承的文献层与不继承的层}
van Boekel 说明食品品质动力学用速率方程，Arrhenius 只是其中一种温度依赖\cite{vanboekel2008}。本文继承速率方程这一层，温度依赖在规格中仍用 Arrhenius，并保留其适用范围：温度窗口、水分开关和反应物是否耗尽都要写明。不把所有品质都写成一级反应。

Holdsworth 关于热力杀菌的标准表述把 $F$ 值沿冷点历史积分\cite{holdsworth1997}。本文把它写成进度 $\alpha_m$，不把 $\alpha_{\mathrm{safe}}$ 写成法规限值。$\alpha_{\mathrm{safe}}$ 与分位、滞回一样是设计量。

蛋白变性、淀粉糊化和美拉德的动力学实验表明三类进度对水分的开关不同。本文写成 $g_i(\varepsilon_l,S)$。不继承用颜色代替中心进度。美拉德文献只解释表面色泽，也就是路径 B 所接触的那一类表面反应。它不提供冷点是否达到热力进度的证据。微生物文献只提供积分形式。初始污染量不在状态向量中，绝对菌数也不进入决策。

\paragraph{方程、算法或协议}
规格中的速率是
\[
\dot\alpha_i=A_i\exp\big(-E_{a,i}/RT\big)\,g_i(\varepsilon_l,S)\,\alpha_i^{n_i},
\]
其中 $n_i$ 不必为 $1$。热力当量沿冷点写成
\[
F=\int_0^t 10^{(T_{\mathrm{cold}}(s)-T_{\mathrm{ref}})/z}\,\mathrm{d}s,
\]
再映射到 $\alpha_m$。评分核没有使用 Arrhenius 形式，也没有分别积分蛋白、淀粉和美拉德。它只在冷点上累积
\[
\dot f=10^{(T_{\mathrm{core}}-70)/10}/60,\qquad\alpha_m=1-\mathrm{e}^{-f/2}.
\]
参照温度 $70^\circ\mathrm{C}$，$z=10\,\mathrm{K}$，时间以秒计再除以 $60$，所以 $f$ 的单位是分钟量级的等效时间。映射 $\alpha_m=1-\mathrm{e}^{-f/2}$ 把等效时间压进 $(0,1)$。当 $f\to\infty$ 时 $\alpha_m\to 1$，于是 $e_\alpha=\alpha_m-0.85$ 的上沿是 $+0.15$。这个上沿是饱和，不是测量精度。

表面褐变在评分核里是另一条经验式：表面温度超过 $110^\circ\mathrm{C}$ 且表面含水低于 $0.55$ 时才增加。它不进入 $\alpha_m$。因此同一条温度史上，中心进度和表面色泽可以分开。这与“颜色不代替中心进度”一致，但系数 $2\times 10^{-4}$ 没有文献区间，只能留在实现里，不能进入主结果的参数表。

设计线 $0.85$ 只用于定义 $e_\alpha$。停机时 $\alpha_m<0.85$ 记为负误差，表示停在设计线之下。它不是认证通过线。腔温闭环的中位 $e_\alpha$ 为 $+0.15$，对应进度已经饱和，不能读成“比设计线精确高出 $0.15$”。

\subsection{来源对照}
\begin{table}[ht]
\centering
\caption{进度，而不是温度阈值。}
\begin{tabularx}{\textwidth}{lXll}
\toprule
来源 & 使用的结果 & 方式 & 不继承 \\
\midrule
van Boekel, 2008 & 品质用速率方程 & 速率层 & 全部写成一级反应 \\
Holdsworth, 1997 & \(F\) 沿冷点历史积分 & \(\alpha_m\) 的形式 & 法规限值 \\
蛋白、淀粉、美拉德传统 & 水分开关不同 & \(g_i(\varepsilon_l,S)\) & 颜色代替中心进度 \\
\bottomrule
\end{tabularx}
\end{table}
后一行没有单独条目。评分核也没有分开这三类积分。

\paragraph{本节不能推出的结论}
$F$ 的积分形式不能推出初始菌数，也不能推出法规限值已经被满足。美拉德路径上的表面色泽不能推出冷点安全。评分核的 $z$ 值映射不能当成 Arrhenius 参数已经拟合。

\section{逆热传导与可观测性}

\paragraph{要判定的量}
表面温度的采样加密以后，内部热阻的信息是否增加。要判定的是线性化秩，以及集合更新里哪些量可以当决策量。稳态、无相变、无内部梯度时，答案应为零灵敏度。

\paragraph{继承的文献层与不继承的层}
Özisik 与 Orlande 给出逆热传导对高频噪声指数放大的结构\cite{ozisik2000}。本文继承病态性：把表面温度历史反推内部热流或内部热阻时，高频成分被放大，正则化是估计问题的一部分。不把他们的某一正则化算法当成唯一估计器。

Beck 的序贯估计传统用未来时间步稳定热流估计。本文只把它当作病态程度的参照：需要向未来看多远，说明当前观测有多弱。不作为在线主算法，也不把整机问题降成单层壁面。

线性系统的可观测性与 Kalman 更新提供秩条件和协方差更新。本文只在两节点命题里使用这一层。不继承“线性模型覆盖相变”。在 \eqref{eq:twonode} 里潜热与腔温的输入列共线；离开平衡之后，气流段仍可以使 \(R_{\mathrm{ext}}\) 与 \(R_{\mathrm{int}}\) 的灵敏度列分开。线性化只在相变方向已经给定的短区间内使用。

Evensen 给出集合卡尔曼滤波的预测—分析格式\cite{evensen2003}。本文把它定为场层估计的规格：集合先按模型推进，再用观测做分析更新。不把集合均值当成决策量。停止看的是进度的下分位。集合均值“看起来已经熟了”不能通过。

\paragraph{方程、算法或协议}
两节点线性化后，仅腔温扰动时输入矩阵是单列。潜热也只进入表面节点，输入列与腔温列共线，变号不增加秩。紧接着只改外部热阻的气流段，使 $(R_{\mathrm{ext}},R_{\mathrm{int}})$ 的灵敏度列不能互为标量倍。这是本文在蒸烤边界上的推论，不是\cite{ozisik2000}或\cite{evensen2003}的原定理。证明自洽地比较两段特解：冷凝项进入表面节点时不经过外部热阻；气流段只改外部热阻。输入列共线并不取消这个参数列的结论。

数值核对用的不是逆问题算法。两段脉冲的表面—内部滞后为 $40.846$，单独蒸汽为 $38.912$，判据是差值大于 $1$。没有声称这是 $1.15$ 倍。稳态极差为 $0$，表示平衡温度里读不出三组热阻的差别。

评分核的在线估计不是集合卡尔曼滤波。它是标量滞后
\[
\widehat T_c\leftarrow\widehat T_c+\Delta t\,(y_T-\widehat T_c)/\tau,
\]
再把估计压到表面观测之下，并用探针窗口更新一个标量 $\hat r$ 及其方差。方差的先验是 $0.64$。更新被接受时乘 $0.35$，即降到 $0.224$。这不是 $\log\det$ 的信息增益，也没有集合。集合规模 $32$ 与 $128$ 没有进入主日志。在线 Fisher 矩阵没有运行。

噪声幅度 $0.4^\circ\mathrm{C}$ 是确定性散列，不是\cite{ozisik2000}里用来展示指数放大的随机高频序列。用该幅度不能回推逆问题的放大倍数。

\subsection{来源对照}
\begin{table}[ht]
\centering
\caption{表面温度为什么不够。序贯估计传统没有单独条目。}
\begin{tabularx}{\textwidth}{lXll}
\toprule
来源 & 使用的结果 & 方式 & 不继承 \\
\midrule
Özisik、Orlande, 2000 & 高频噪声被放大 & 病态结构 & 唯一正则化算法 \\
序贯估计传统 & 未来步稳定热流估计 & 病态参照 & 整机降成单层壁 \\
可观测性与 Kalman 更新 & 秩与协方差 & 两节点命题 & 线性模型覆盖相变 \\
Evensen, 2003 & 预测—分析格式 & 场层规格 & 集合均值当决策量 \\
\bottomrule
\end{tabularx}
\end{table}

\paragraph{本节不能推出的结论}
稳态极差为 $0$ 只支持“平衡观测不含内部热阻”，不支持内部状态已经被重构。标量方差下降不支持集合卡尔曼滤波已经运行。病态性文献不提供本文的正则化参数。

\section{信息价值与对偶控制}

\paragraph{要判定的量}
一次阀门动作是在推动品质，还是在做一次实验。要判定的是信息有没有下降；没有下降时，停止时刻是否保持不动。家电执行器不能发出任意波形，因此还要判定探针库是不是该问题的可执行子集。

\paragraph{继承的文献层与不继承的层}
Lindley 把实验提供的信息定义为先验熵与后验熵之差\cite{lindley1956}。本文的 $G(a)$ 采用这一层。不把效用只写成参数熵。决策熵看的是 $\vect{z}$ 里与停止、表面品质和均匀性有关的分量，不是全部物性参数的熵。

Fedorov、Kiefer 等最优实验设计给出 $\log\det$ 与参数置信椭球。本文只把 Fisher 形式当作对照。不把离线一次性设计代替在线脉冲。二者排序不一致时以集合形式为准；而集合形式本身尚未作为主结果跑完，所以 Fisher 对照也还没有排序可报。

Feldbaum 的对偶控制说明控制量同时影响状态与信息\cite{feldbaum1965}。本文继承问题提法：输入既出现在状态方程里，也出现在信息增益里。不继承其原问题的解析解。家电的加热、蒸汽、风速和排汽不是任意波形。

谨慎控制与主动学习控制把有限时域上的信息与代价并列。探针库是该问题的可实现子集。不把家电执行器当成任意波形发生器。

\paragraph{方程、算法或协议}
信息增益的规格是
\[
G(a)=\frac12\log\det\Sigma_z(t)-\frac12\log\det\Sigma_z(t+\tau;a).
\]
代价是能量、失水、表面褐变和时间的加权。探针只在 $G$ 超过下限且损伤低于预算时执行。P4 的顺序由秩条件固定：先相变变号，再对流切换。无效则停止时刻不动。评分核没有计算 $G(a)$。它用固定窗口 $[150,190)\,\mathrm{s}$ 代替在线打分，用“是否更接近 $r_{\mathrm{proxy}}$”代替熵差，用方差乘 $0.35$ 代替 $\log\det$ 的下降比例。窗口不随赤字开关。赤字大于零但方差已经不大于 $0.25$ 时，核不再进入探针。这是一条阈值规则，不是在探针库上最大化信息增益减去代价。

厚件上还有一条与信息有关、但不是熵的否决：对偶且方差仍大于 $0.45$、厚度大于 $0.04\,\mathrm{m}$ 时，即使冷点估计已经越过裕量，也不停止。这是否决，不是 $G(a)$ 的数值。

厚度上的定性只能引用第 10 章已经汇总的交叉：薄件上探针没有把对偶从无探针优化里分开，两者都停在设计线之下；中等厚度上对偶把中位误差从约 $-0.53$ 拉到约 $-0.08$；厚件上两者都提前，对偶只略好。精确交叉格以第 10 章的表为准。这些差别是停止误差，不是 $G(a)$ 的测量。

\subsection{来源对照}
\begin{table}[ht]
\centering
\caption{动作为什么可以是一次实验。最优设计与谨慎控制没有单独条目。}
\begin{tabularx}{\textwidth}{lXll}
\toprule
来源 & 使用的结果 & 方式 & 不继承 \\
\midrule
Lindley, 1956 & 信息是熵差 & \(G(a)\) 的层 & 效用只写参数熵 \\
最优设计传统 & \(\log\det\) 与置信椭球 & Fisher 对照 & 离线设计代替脉冲 \\
Feldbaum, 1965 & 输入同时推动状态与信息 & 问题提法 & 原问题的解析解 \\
谨慎控制传统 & 有限时域并列信息与代价 & 探针库 & 任意波形 \\
\bottomrule
\end{tabularx}
\end{table}

\paragraph{本节不能推出的结论}
对偶问题的提法不能推出解析最优波形已经得到。固定窗口不能写成在线信息最大化。方差乘 $0.35$ 不能写成 Lindley 熵差。中等厚度上的误差变化不能单独证明探针库优于连续对偶控制。

\section{算子学习}

\paragraph{要判定的量}
神经算子是加速器还是定律。要判定的是同一训练集上，哪一个候选在蒸汽干预的位移误差上更低。文献不提供这个胜负。

\paragraph{继承的文献层与不继承的层}
Raissi、Perdikaris 与 Karniadakis 给出用残差惩罚方程的训练方式\cite{raissi2019}。本文只把它当作训练时的物理损失项。不用该网络代替在线安全判断。安全投影的自变量不含网络输出。

Li 等给出用 Fourier 神经算子学习参数化演化的结构\cite{li2021fno}。本文把它列为候选 $\mathcal{G}_\phi$。不继承训练分布之外的精度保证。

Lu 等给出 DeepONet 的分支—主干结构\cite{lu2021deeponet}。本文把它列为同一数据上的对照算子。不继承“一种结构优于另一种”的普遍结论。

\paragraph{方程、算法或协议}
候选算子学习的是边界到决策状态的映射，而不是一个新的本构关系。选择规则是：同一训练集上，蒸汽干预造成的位移误差更低者为主模型。另一个保留为对照。残差超过预设上界时回退到物理步进，不把算子输出送进屏障。

Fourier 算子和 DeepONet 仍停留在规格，二者之间没有位移误差的表，也没有残差界的回退次数。另有一次训练写在演化一章：教师是两节点滞后，结构是主干加一支同形分支，留出均方误差 $1683.27\,(^\circ\mathrm{C})^2$，只留主干时 $2106.23\,(^\circ\mathrm{C})^2$。这个差不是两种文献结构的胜负。评分核的十六条步进仍是战役日志里的演化。它的截断误差来自空间步长 $L/16$ 和固定步长 $2\,\mathrm{s}$。

主模型尚未在 Fourier 算子与 DeepONet 之间选出，因此任何“算子已经加速了试验”的句子都没有对象。教师扩充标记也没有进入主表。

\subsection{来源对照}
\begin{table}[ht]
\centering
\caption{加速器，不是定律。}
\begin{tabularx}{\textwidth}{lXll}
\toprule
来源 & 使用的结果 & 方式 & 不继承 \\
\midrule
Raissi 等, 2019 & 残差惩罚方程 & 训练损失项 & 网络代替安全判断 \\
Li 等, 2021 & Fourier 算子 & 候选 \(\mathcal{G}_\phi\) & 训练外的精度保证 \\
Lu 等, 2021 & 分支—主干 & 同一数据上的对照 & 一种结构普遍更优 \\
\bottomrule
\end{tabularx}
\end{table}

\paragraph{本节不能推出的结论}
两篇算子文献不能推出谁在蒸烤干预上更准。物理损失项不能推出在线安全已经由网络保证。未运行的回退次数不能画成曲线。

\section{因果、预测控制与安全过滤}

\paragraph{要判定的量}
蒸汽干预的效应能否按路径相加，以及安全可行域是否随后验变形。要判定的还有：硬回路的存在是算法结论，还是器具安全框架提出的要求。

\paragraph{继承的文献层与不继承的层}
Pearl 区分预测、干预和反事实，并把 $\mathrm{do}$ 定义成替换结构方程\cite{pearl2009}。本文按这个顺序计算路径 A、B、C：先预测，再替换蒸汽节点，最后才在需要时做反事实对比。不从短烹饪序列里发现父节点。父节点是规格里写死的，不由数据搜索改写。

Rawlings、Mayne 与 Diehl 给出有限时域约束优化这一类问题\cite{rawlings2017}。本文的预测控制属于这一类。不继承某一求解器的默认最优性声明。

Ames、Xu、Grizzle 与 Tabuada 给出屏障函数约束下的二次投影\cite{ames2017cbf}。本文继承投影形式：在安全约束下，把动作改到离参考动作最近的可行点。不继承“屏障条件自动等于器具安全认证”。

IEC 60335-1 与 GB 4706.1 给出家用电器在漏电、温升和异常工作上的安全框架\cite{iec60335,gb4706}。本文只继承“硬回路必须存在，且不读取品质后验”这一要求来源。不把标准条款写成算法已经满足。实施时用的是现行条款的入口，不是本文重新编号的合格声明。

\paragraph{方程、算法或协议}
路径可加是链式法则。少掉 C，结壳后期没有项。短探针窗口内 C 是小量，全时域积分必须保留。干预计算把蒸汽节点替换为接管值；条件相关则保留策略到蒸汽的边。二者不是同一个数。

预测控制的规格是有限时域上最小化品质代价，约束写成机会约束，探针占用区间内输入等于已经发布的探针。求解器被指定为 CasADi/IPOPT。这次没有调用该求解器。评分核里名为无探针优化和对偶的两支，用的是估计温度的饱和线性请求，不是有限时域优化的解。因此战役日志不能支持“求解器已声明最优”。控制一章另有一次 18 点搜索：一个状态、60 s、无约束，得分最低的加热是 0.7，饱和律给出 1.0。那一次搜索同样没有声明网格之外的最优。

投影的规格是
\[
\vect{u}=\argmin_{\vect{u}}\|\vect{u}-\vect{u}_{\mathrm{AI}}\|_W^2
\]
并满足 $\dot h_j+\gamma_j h_j\ge 0$。$h_j$ 的自变量只有漏电、水位或干烧、加热器温度、腔压和门。后验只通过参考点进入目标。命题是：可行域与后验不交。断言实现的是一条测试函数，不是二次规划：加热器超过 $210^\circ\mathrm{C}$ 时把请求剪到 $0.25$，两个不同后验得到同一个剪裁；加热器低于该阈值时保持 $1$。物理核的削减从 $230^\circ\mathrm{C}$ 开始，按 $(T_h-230)/40$ 最多削减 $75\%$，作用在已经形成的请求上。$210$ 与 $230$ 是两个对象。测试函数通过，只说明“不随后验改变”这一方向；它不说明二次规划已经求解，也不说明标准条款已经满足。

硬件切断不读估计器。这一句是结构要求。本次日志没有漏电、水位、腔压和门的采样，不能用加热器削减代替整条硬回路。

\subsection{来源对照}
\begin{table}[ht]
\centering
\caption{解释与执行。}
\begin{tabularx}{\textwidth}{lXll}
\toprule
来源 & 使用的结果 & 方式 & 不继承 \\
\midrule
Pearl, 2009 & \(\mathrm{do}\) 是替换结构方程 & 路径顺序 & 从短序列发现父节点 \\
Rawlings 等, 2017 & 有限时域约束优化 & 问题类 & 求解器默认最优 \\
Ames 等, 2017 & 约束下的二次投影 & 投影形式 & 自动等于器具认证 \\
IEC 60335-1，GB 4706.1 & 漏电、温升、异常工作 & 硬回路必须存在 & 条款写成算法已满足 \\
\bottomrule
\end{tabularx}
\end{table}

\paragraph{本节不能推出的结论}
路径公式不能从短序列里发现新的父节点。饱和线性请求不能写成模型预测控制的最优解。屏障测试函数通过不能写成器具已经按安全标准认证。标准文本不提供软件投影的数值参数。

\section{创新所在的层}

\paragraph{要判定的量}
新的对象是一条方程，还是几条已经存在的方程之间的连接。本节把前面各节的“不继承”列收束成四条连接，并规定怎样退出创新表述。

\paragraph{继承的文献层与不继承的层}
多孔介质方程、熵差、$\mathrm{do}$ 算子、神经算子和屏障函数都已经有来源\cite{whitaker1977,lindley1956,pearl2009,li2021fno,lu2021deeponet,ames2017cbf}。这些层全部不作为本文新创造的对象。不继承“把连接说成新的本构、新的熵或新的屏障定义”。

\paragraph{方程、算法或协议}
新增的是四条连接。

\begin{enumerate}
\item 决策状态 $\vect{z}$ 从全场中取出，并与稳态不可辨识接在一起：平衡且无相变时，内部热阻不进入代数方程，灵敏度列消失。稳态腔温再密也不更新 $R_{\mathrm{int}}$。信息动作必须离开平衡。
\item 信息动作落成 P4：相变变号之后接对流切换。方差或分位没有下降时，停止时刻保持在动作之前。这是无效探针否决，不是新的熵定义。
\item 同一次蒸汽干预分解为三条路径：潜热到冷点，表面水分到品质，结壳到内部热阻再到冷点。少一条，结壳后期的热阻没有解释项。
\item 品质后验与安全可行域使用不相交的测量。品质链可以错，可行域的形状不随后验改变。硬件切断不读估计器。
\end{enumerate}

这四条若与某篇已发表论文的技术方案相同，改为引用并退出创新表述。本文没有把“未读到”写成“不存在”。肉类等温线、公开的烤制中心温度历程、以及与这四条连接完全相同的技术方案，都只在已经读到的范围内讨论。

评分核实现了这四条里的一部分，而且是简化以后的部分。决策状态在日志里只保留冷点、表面、进度、热阻代理和方差，没有表面褐变进入停止规则，没有温度离散的决策项。P4 被固定时间窗代替，否决只在厚件且方差仍高时出现。三条路径没有在评分核里按链式法则计算，浏览器里的三条显示公式更不能算实现。测量隔离只落实到加热器温度这一路硬测量；漏电、水位、腔压和门没有进入这次步进。因此四条连接是本文的结构主张，不是这次日志已经逐项跑完的清单。

\paragraph{本节不能推出的结论}
四条连接不能改写 Fourier 定律，不能替换熵的定义，也不能替换屏障函数的定义。部分实现不能写成四条都已在主结果里定量成立。未检索到相同技术方案，只说明本次阅读没有遇到，不说明文献中没有。
